%% -*- coding: utf-8 -*-
\documentclass[12pt,pagesize,paper=landscape,paper=192mm:108mm]{scrbook} 
%1920x1080 1280x720
\areaset[current]{192mm}{108mm}
\usepackage{calc}
\usepackage[T2A]{fontenc}
\usepackage[utf8]{inputenc}
\usepackage[english,russian]{babel}
\usepackage{microtype}
\usepackage{misccorr}
\usepackage{cmap}
%\usepackage[unicode=true]{hyperref}
\usepackage{graphicx}
\usepackage{amssymb}
\usepackage{amsmath}
%\usepackage{srcltx}
\usepackage{textcomp}
\usepackage{braket}
\usepackage{xspace}
%научные символы и смайлики \smiley \frownie
\usepackage{wasysym}
\usepackage{ccicons}
\DeclareMathOperator{\Tr}{Tr}
\DeclareMathOperator{\Texp}{Texp}
\renewcommand{\epsilon}{\varepsilon}
%перенос формул в тексте
\newcommand*{\hm}[1]{#1\nobreak\discretionary{}%
  {\hbox{$\mathsurround=0pt #1$}}{}}
\begin{document}
\begin{titlepage}
  \vspace*{-0.5em}
  \begin{center}    
    \hspace*{3em}
    \begin{minipage}[t]{3em}
      \includegraphics[width=\textwidth]{../BINP-logo}
    \end{minipage}\hfill
    \begin{minipage}{0.23\linewidth}
    \includegraphics[width=\textwidth]{../NSU-logo}
    \end{minipage}
    \hfill
    \hspace*{6em}

    Кафедра теоретической физики физического факультета НГУ
    \medskip

    \Large
    Резниченко А.\,В.
    \bigskip

    \huge
    \textbf{Квантовая электродинамика}
    \bigskip

    \Large
    Семинар № 3
    \vfill

    \normalsize
    \vfill

    \normalsize \ccbysa\hspace{0.5em}  Новосибирск 2015
  \end{center}
\end{titlepage}
\newpage

\vspace*{-1em}
\begin{center}
\vfill
  \begin{minipage}{0.7\linewidth}

    Пропагаторы фотона и электрона (продолжение).  Задача: показать,
    основываясь на канонических коммутационных соотношениях, что
    пропагатор фотона $D^{\mu\nu}(x)$ в ковариантной калибровке с
    $\xi=1$ является функцией Грина оператора $\partial^2$:
    $\partial^2 D^{\mu\nu}(x)=i g^{\mu\nu}\delta(x)$.  Задача:
    показать, основываясь на антикоммутационных соотношениях, что
    $\delta(x^0-y^0)[j^0(x),j^{\mu}(y)]=0$, где
    $j^{\mu}={:}\bar{\Psi}(x) \gamma^{\mu} \Psi(x){:}$ "--- оператор
    электромагнитного тока.
    \smallskip

    Унитарность $S$"=матрицы в ковариантной калибровке.  Физическое
    пространство фотонных состояний $\ket{\Phi}$:
    $\partial_{\mu}A^{\mu+}\ket{\Phi}=0$, где $A^{\mu+}$ "---
    положительночастотная часть полевого оператора
    $A^{\mu}$. Следствия соотношения $[\partial_{\mu}A^{\mu+}, S]=0$:
    физические состояния переходят в физические состояния, оператор
    эволюции $S=\Texp[i \int L_{\text{int}} d^4 x]$ унитарен на
    физическом пространстве.
  \end{minipage}
  \vfill

  % \normalsize \ccbysa\hspace{0.5em} Новосибирск 2013
\end{center}
\end{document}
